\chapter{بستارها و عبارت‌های لامبدا}%
\label{ch:lambdas}

تا اینجا روال‌ها را همیشه با نام می‌شناختیم. اما در سلام، روال خودش یک \emph{مقدار}
است: می‌توان آن را در متغیر گذاشت، به روالِ دیگر داد، یا از یک روال
بازگرداند. در این فصل با \textbf{روال‌های بی‌نام} (لامبدا) و \textbf{بستارها}
آشنا می‌شویم ابزارهایی که برنامه‌نویسی را انعطاف‌پذیرتر می‌کنند.

\section{روال به‌مثابهِ مقدار}

به این مثالِ واقعی نگاه کنید. اینجا روال‌ها را در متغیر می‌گذاریم و حتی به روالِ
دیگری می‌دهیم:

\begin{salamfa}
روال اعمال(ف : روال (صحیح) صحیح, مقدار : صحیح) : صحیح:
    برگشت ف(مقدار)
پایان

روال ریشه:
    دوبرابر : روال (صحیح) صحیح = (ن : صحیح) => ن * 2
    مربع : روال (صحیح) صحیح = (ن : صحیح) => ن * ن
    سرچاپ اعمال(دوبرابر, 21)
    سرچاپ اعمال(مربع, 9)
    سرچاپ دوبرابر(مربع(3))
پایان
\end{salamfa}

چند چیزِ تازه اینجا هست:

\subsection{گونهِ روال}

\begin{salamfa}
روال (صحیح: ) : صحیح
\end{salamfa}

این یک \textbf{گونه} است مانندِ \code{صحیح} یا \code{رشته}، اما گونهِ یک
\emph{روال}. خوانده می‌شود: «روالی که یک عددِ صحیح می‌گیرد و یک عددِ صحیح
برمی‌گرداند.» پس \code{دوبرابر} متغیری است که نه عدد، بلکه یک \emph{روال} در
آن نگه‌داشته می‌شود.

\subsection{عبارتِ لامبدا}

\begin{salamfa}
(ن صحیح) => ن * 2
\end{salamfa}

این یک \textbf{لامبدا} است: یک روالِ کوچکِ بی‌نام. ساختارش ساده است: درونِ
پرانتز پارامترها، سپس نشانهٔ \code{=>}، و سپس عبارتی که نتیجه را می‌دهد. این
لامبدا یک عدد می‌گیرد و دوبرابرش را برمی‌گرداند درست مانندِ یک روالِ معمولی،
اما بدونِ نام و در یک خط.

\subsection{روالی که روال می‌گیرد}

روالِ \code{اعمال} در پارامترِ نخستش یک \emph{روال} می‌گیرد (\code{ف}) و آن را
روی \code{مقدار} اجرا می‌کند. به همین دلیل توانستیم هم \code{دوبرابر} و هم
\code{مربع} را به آن بدهیم. روال‌های که روالِ دیگر می‌گیرند یا برمی‌گردانند را
\textbf{روال‌های مرتبهٔ بالاتر} می‌نامند و در پردازشِ داده بسیار کاربرد دارند.

\begin{note}
به آخرین خط دقت کنید: \code{دوبرابر(مربع(3))}. نخست \code{مربع(3)} برابرِ ۹، سپس
\code{دوبرابر(9)} برابرِ ۱۸ می‌شود. چون این‌ها متغیرهایی حاوی روال‌اند، می‌توان
مستقیم صدایشان زد درست مانندِ روال‌های معمولی.
\end{note}

\section{بستار: لامبدایی که محیطش را به خاطر می‌سپارد}

اینجا به یکی از قدرتمندترین ایده‌های برنامه‌نویسی می‌رسیم. یک لامبدا می‌تواند
متغیرهای \emph{بیرونِ} خودش را «به یاد بسپارد». به چنین لامبدایی \textbf{بستار}
(closure) می‌گویند. این مثالِ واقعی را ببینید:

\begin{salamfa}
روال سازنده جمع(ن : صحیح) : روال (صحیح) صحیح:
    برگشت (x : صحیح) => x + ن
پایان

روال ریشه:
    دوبرابر : روال (صحیح) صحیح = (n : صحیح) => n * 2
    سرچاپ دوبرابر(21)
    جمع۱۰ : خودکار = سازنده جمع(10)
    سرچاپ جمع۱۰(5)
پایان
\end{salamfa}

روالِ \code{سازنده جمع} یک عدد (\code{ن}) می‌گیرد و یک \emph{لامبدا}
برمی‌گرداند. اما دقت کنید: آن لامبدا \code{(x صحیح) => x + ن} از
\code{ن} استفاده می‌کند که پارامترِ روالِ بیرونی است! وقتی
\code{سازنده جمع(10)} را صدا می‌زنیم، لامبدای بازگشتی مقدارِ \code{ن = 10} را
\emph{با خود نگه می‌دارد}. پس \code{جمع۱۰(5)} برابرِ $5 + 10 = 15$ می‌شود.

این یعنی \code{جمع۱۰} یک روالِ تازه است که «۱۰ را به یاد دارد». اگر
\code{سازنده جمع(100)} می‌ساختیم، روالی می‌داشتیم که ۱۰۰ را به یاد می‌آورد.

\section{بستار با چند دستور}

لامبداهایی که دیدیم یک عبارت داشتند. اما گاهی به چند خط منطق نیاز داریم. برای
این، شکلِ بلندترِ لامبدا را داریم که با \code{=> :} باز و با \code{پایان} بسته
می‌شود. این مثالِ واقعی یک \emph{شمارنده} می‌سازد:

\begin{salamfa}
روال ساخت شمارنده() : روال () صحیح:
    ناپایا ن : صحیح = 0
    برگشت () => :
        ن = ن + 1
        برگشت ن
    پایان
پایان

روال ریشه:
    بعدی : روال () صحیح = ساخت شمارنده()
    سرچاپ بعدی(), بعدی(), بعدی(), بعدی()
پایان
\end{salamfa}

اینجا اتفاقِ جالبی می‌افتد. \code{ساخت شمارنده} متغیرِ \code{ن} را می‌سازد و
یک بستار برمی‌گرداند که آن را یکی زیاد و سپس بازمی‌گرداند. هر بار که
\code{بعدی()} را صدا می‌زنیم، همان \code{ن} که \emph{میانِ فراخوانی‌ها زنده
می‌ماند} یکی بیشتر می‌شود. پس خروجی \code{1 2 3 4} است.

\begin{deep}[نگاهی به درون: بستار حافظه را کجا نگه می‌دارد؟]
وقتی یک بستار متغیری از محیطِ بیرون را می‌گیرد، سلام آن متغیر را در \emph{پشته}
نمی‌گذارد (که با پایانِ روال از میان می‌رود)، بلکه در \emph{پشتهٔ آزاد} (heap)
نگه می‌دارد و یک «محیطِ بسته‌شده» می‌سازد که همراهِ خودِ روال زندگی می‌کند. به
همین دلیل است که \code{ن} حتی پس از پایان‌یافتنِ \code{ساخت شمارنده} همچنان
زنده می‌ماند. این مکانیزم گرفتنِ متغیرها و نگه‌داشتنشان روی پشتهٔ آزاد ---
قلبِ کارِ بستارهاست.
\end{deep}

\begin{tip}
بستارها به‌ویژه برای ساختنِ «کارخانه‌های روال» (مانندِ \code{سازنده جمع}) و
نگه‌داشتنِ حالتِ خصوصی (مانندِ شمارنده) عالی‌اند. در فصلِ کتابخانهٔ استاندارد
خواهیم دید که بسیاری از روال‌های پردازشِ داده، لامبدا می‌گیرند تا رفتارشان را
سفارشی کنید.
\end{tip}

\section{خلاصهٔ فصل}

\begin{itemize}
  \item در سلام، روال یک مقدار است: می‌توان آن را در متغیر گذاشت، پاس داد و
        بازگرداند.
  \item گونهِ روال چنین نوشته می‌شود: \code{روال (ورودی) خروجی}.
  \item لامبدا روالی بی‌نام است: \code{(پارامتر) => عبارت}.
  \item برای چند دستور، شکلِ بلند را به‌کار می‌بریم: \code{() => :} \ldots
        \code{پایان}.
  \item بستار لامبدایی است که متغیرهای محیطِ خود را به یاد می‌سپارد.
  \item روال‌های مرتبهٔ بالاتر، روالِ دیگر را می‌گیرند یا برمی‌گردانند.
\end{itemize}

\section{تمرین‌ها}

\exercise یک لامبدا بسازید که عددی بگیرد و مکعبش را برگرداند، و آن را در یک
متغیر بگذارید و صدا بزنید.

\exercise روالی به نامِ \code{سازنده ضرب} بنویسید که عددی بگیرد و یک بستار
برگرداند که هر ورودی را در آن عدد ضرب کند.

\exercise با کمکِ \code{ساخت شمارنده}، دو شمارندهٔ مستقل بسازید و نشان دهید که
هر کدام حالتِ خودش را دارد.

\exercise روالِ \code{اعمال دوبار} را بنویسید که یک روال و یک مقدار بگیرد و
روال را \emph{دو بار} پشتِ سرِ هم روی مقدار اجرا کند.

\exercise بستاری بسازید که مجموعِ همهٔ اعدادی را که تا کنون به آن داده‌اید نگه
دارد و در هر فراخوانی، مجموعِ تازه را برگرداند.
